\chapter{Background}
\label{ch:background}

Five bodies of work bear on this system: worn robot limbs and what they are
for, the split between the person wearing a machine and the person driving it,
the problem of driving a robot from a control that is not the same shape,
sharing control between a person and a program, and how the wearer's
experience is measured. This chapter takes them in that order and finishes
with what is left open, which is narrower than it first looks.

\section{Robot limbs worn on the body}

The field's founding demonstrations come from Parietti and Asada, whose extra
limbs brace against the environment to steady the wearer while they work
overhead~\cite{parietti2014bracing,parietti2013dynamic}, later extended to
supporting the whole body~\cite{parietti2016support}. Bracing is a useful
starting point because it names the mechanical difficulty at once. An arm
mounted on a person pushes back into that person, and the person moves.

Later reviews map the design space along consistent
lines~\cite{yang2021srlreview,prattichizzo2021augmentation,tong2021srlreview,zhang2023srlreview}.
Where the limb is mounted trades reach against the twisting load it puts on
the wearer. How it is actuated trades responsiveness against weight. How it is
commanded ranges from muscle signals to a deliberate manual control. A recent
survey adds a pointed observation: the field reports how embodied the wearer
feels far more often than it reports whether the task got
done~\cite{zheng2026quantitative}.

Weight binds harder than anything else and is discussed least. Each Gen3 arm
is \SI{8.2}{\kilo\gram}, so the pair with the harness passes
\SI{17}{\kilo\gram}. Simulation of the loads such systems place on the body
suggests this is not a detail~\cite{moon2024impact}, and it is the main reason
the session design in \cref{ch:study} caps how long anybody wears the rig.

\begin{figure}[htbp]
  \centering
  \reproduced{42mm}{Parietti and Asada's bracing arrangement, with the limbs
  anchored against the work surface so that the wearer, the limbs and the
  environment form one closed chain. Source: \cite{parietti2014bracing}.}
  \caption[Bracing limbs (to be reproduced)]{The arrangement that set the
  mechanical framing for the field. To be reproduced from
  \cite{parietti2014bracing} once permission is obtained.}
  \label{fig:repro-bracing}
\end{figure}

\section{When the wearer and the driver are different people}

The closest predecessor to the arrangement used here is Fusion, which put a
robot head and two arms on a worn frame and let a remote operator drive
them~\cite{saraiji2018fusion}. Two qualifications have to travel with that
reference. Its three modes describe how firmly the robot arms are strapped to
the wearer's own arms, not how much the machine decides; in all three the
machine decides nothing. And it is a two-page demonstration with no user study
and no performance figures. It shows the arrangement existed. It cannot
support a claim about how well it worked.

The evidence about how it feels to wear one of these comes from elsewhere.
Recent work on how close such a limb should come to the body ran a hidden
human operator, varied how autonomous the machine appeared to be, and measured
the wearer's arousal and reported trust~\cite{zhou2026proxemics}. The result
runs against intuition. Greater apparent autonomy made both worse: arousal up,
trust down. Because the autonomy in that study was a concealed person and the
wearer had nothing else to do, it does not settle what happens when the
autonomy is real and the wearer is busy. It does establish that the wearer's
experience cannot be assumed benign.

Part of the two-person problem is already solved. Zhang et al.\ compensate for
the motion a human base introduces and report end-effector errors around
\SI{1.4}{\milli\metre} with the shoulder as a moving
base~\cite{zhang2024motioncomp}. What that leaves open is the case here, where
the disturbance comes from a nervous system the driver has no advance warning
of. When one person both wears and drives, the sway is partly predictable to
the controller because the same person caused it. With two people it is not.

\section{Driving something that is not the same shape}

If the control and the robot differ in size or proportion, the mapping between
them has to be chosen. Guggenheim and Asada argue for small physical copies of
the robot on the grounds that a similar shape gives the operator feedback for
free, without a separate force display~\cite{guggenheim2021haptic}. The
argument holds for similarity of shape and says nothing about similarity of
size, which is where this system departs: the master's links total
\SI{0.272}{\metre} against the robot's \SI{0.902}{\metre} of reach, so the
same joint angles produce very different hand movements.

Two mappings are available. Sending master joint angles straight to robot
joint angles preserves the shape the arm takes, which matters when the shape
itself has to avoid something, but it demands that every master joint be
measured accurately. Mapping hand movement to hand movement decouples the two
chains and leaves the robot's spare joint unspecified. \Cref{ch:master}
records why this project started with the first and ended with the second.

A related question is what to feed back to the driver. Pinardi et al.\ review
the added sensory channels used in movement augmentation and find that
improving control and improving the sense of ownership do not go
together~\cite{pinardi2023feedback}. Neither is a proxy for the other, which is
why \cref{ch:study} measures both.

\section{Sharing control with the machine}

Shared autonomy asks a different question from teleoperation: not how to
transmit the operator's intent faithfully, but how much of it the machine
should supply. Policy blending mixes the operator's command with the
machine's estimate of the goal in proportion to
confidence~\cite{dragan2013policyblending}. Hindsight optimisation treats it
as a decision under uncertainty about which goal is
meant~\cite{javdani2015hindsight}. Both assume a bench robot and an operator
who can see it directly.

Two things change when the robot is worn by a second person. The cost of
guessing wrong is no longer a dropped object, and the operator's confidence is
no longer the only confidence that matters. The design used here keeps
position entirely under the driver's control in every mode and lets the
machine supply only the wrist orientation, which is the part the failing
sensor channels cannot deliver anyway. That choice is argued in
\cref{ch:system} and is the reason the assistance can be switched on without
changing who is responsible for where the hand goes.

\section{Measuring the wearer's experience}

Three instruments recur in this literature and are used unchanged here. Task
load is measured with the NASA Task Load Index~\cite{hart1988nasatlx}, physical
exertion with the Borg category-ratio scale~\cite{borg1982perceived}, and the
sense that a limb belongs to you with the embodiment questionnaire developed
from the rubber-hand work~\cite{longo2008embodiment,botvinick1998rubber}.
Difficulty of a pointing task is indexed in the standard
way~\cite{fitts1954information}, which lets tasks of different sizes and
distances be compared.

\section{What is open}

It is tempting to claim that nobody has looked at a two-person worn limb with
real autonomy. That would be too strong. The concealed-operator
study~\cite{zhou2026proxemics} has already varied autonomy and measured the
wearer, and Fusion~\cite{saraiji2018fusion} already built the two-person
arrangement. What has not been done is to measure task performance and wearer
experience together, in the same sessions, with autonomy that is genuinely a
program and a wearer who has a job of their own. That is the gap this study is
designed to fill, and \cref{ch:study} sets out how.
